\section{Notebook Generation}
\label{sec:notebook-generation}

Jupyter notebooks \parencite{kluyver2016jupyter_notebooks} are widely used in ML development because 
they allow code, explanatory text, and visual output to be combined 
in a single document, making them well suited for 
experimentation and exploration. Providing a pre-structured notebook for a 
recommended method gives users a concrete starting point 
without having to write boilerplate code from scratch. 
A template-based approach makes it possible to generate 
such notebooks automatically, with content tailored to 
the selected method and the user's problem description. 

MLguide have the possibility to generate a starter notebook for recommended ML 
methods. This relies on the existence of a pre-written template 
that defines the notebook structure and code.
When a user opens a method details page, by clicking on one of the 
recommended methods, 
a notebook is generated automatically from a template and 
displayed as a code preview. The user can then download the 
notebook or open it directly in Google Colab. This section 
describes how templates are structured, how notebooks are 
generated, and how they are made available to the user.

Notebook generation relies on the following tools:

\begin{itemize}
    \item \textbf{Jinja2} \parencite{jinja2}: A Python 
    templating engine used to render notebook templates 
    with method-specific and user-provided content.
    \item \textbf{jupytext} \parencite{Wouts_Jupytext}: 
    Tool for converting Python code into the 
    \texttt{.ipynb} format used by Jupyter and Google Colab.
    \item \textbf{Google Colab}: A cloud-based notebook 
    environment that allows users to run the generated 
    notebook directly in the browser without local setup.
    \item \textbf{GitHub API}: Used to upload the generated 
    notebook so that Google Colab can load it from a 
    publicly accessible URL.
\end{itemize}

Much of the notebook generation approach, including the use of 
Jinja2 templates and 
the GitHub-based Open in Colab mechanism, is inspired by 
Traingenerator \parencite{rieke2021traingenerator}, which 
is described in Section~\ref{sec:template-based-code-generation}. 
Unlike Traingenerator, where the user selects task, 
framework and more parameters manually, MLguide just selects the template 
automatically based on the recommended method that is chosen by the user. 

\subsection{Template Resolution and Rendering}
\label{sec:template-resolution}

The generation process has three steps. First, the 
template resolver maps the selected method to a Jinja2 
template file. Templates are organised by task family 
and method key, following the path structure: \\
\texttt{templates/notebooks/\{family\}/\{method\_key\}.py.jinja}.

Second, the template is rendered with user-specific 
context, including the problem description, method title, 
and task family. This produces a Python source file 
personalised to the user's input.

Third, the rendered Python source is converted to 
\texttt{.ipynb} format using jupytext. Colab-compatible 
metadata is added at this stage, including kernel 
specification and notebook name, so the file can be 
opened in Google Colab without additional configuration.

\subsection{Colab Export}
\label{sec:colab-export}

The user can open the generated notebook directly in 
Google Colab. The notebook is uploaded to a configured 
GitHub repository via the GitHub API, and a Colab URL 
is constructed from the repository name, branch, and 
file path. The browser opens this URL in a new tab, 
where Colab loads the notebook directly from GitHub.

The whole process of generating the notebook and opening it in Colab is 
described in the sequence diagram in Figure~\ref{fig:notebook-generation-sequence}.

\begin{landscape}
\begin{figure}[H]
    \centering
    \includegraphics[width=\linewidth, height=\textheight, keepaspectratio]{Figures/06-implementation/notebook-generation-sequence.png}
    \caption[Notebook generation sequence]{Sequence diagram of the notebook generation process in MLguide.}
    \label{fig:notebook-generation-sequence}  
\end{figure}
\end{landscape}

\subsection{Template Structure and Extensibility}

The template folder structure matches the 
\texttt{mla:ML\_task\_family} classes defined in the 
ontology extension described in 
Section~\ref{sec:ml-ontology-extension} and listed in 
Table~\ref{tab:ml-task-families}. Template files are 
named after the method's ontology IRI
(the unique identifier used in the knowledge graph), 
which gives accurate mapping between the ontology and the file system.

Many ML methods can be applied to more than one type of 
task, as described in 
Section~\ref{sec:machine-learning-methods} and 
Section~\ref{sec:neural-networks-and-deep-learning}. 
Such methods have a separate template in each relevant 
family folder, since different task types require 
different code structure. The correct template is selected 
by inspecting the task chosen in the users problem description. 
Table~\ref{tab:template-structure} shows examples of 
how methods are distributed across family folders.

\begin{table}[H]
\centering
\caption[Template folder structure]{Example of template 
organisation by task family and method key.}
\label{tab:template-structure}
\begin{tabular}{ll}
\toprule
\textbf{Path} & \textbf{Method} \\
\midrule
\texttt{classification/artificial\_neural\_network.py.jinja} 
    & Artificial neural network \\
\texttt{classification/random\_forest.py.jinja} 
    & Random forest \\
\texttt{classification/support\_vector\_machine.py.jinja} 
    & Support vector machine \\
\texttt{classification/logistic\_regression.py.jinja} 
    & Logistic regression \\
\midrule
\texttt{regression/artificial\_neural\_network.py.jinja} 
    & Artificial neural network \\
\texttt{regression/random\_forest.py.jinja} 
    & Random forest \\
\texttt{regression/support\_vector\_regression.py.jinja} 
    & Support vector regression \\
\texttt{regression/linear\_regression.py.jinja} 
    & Linear regression \\
\midrule
\texttt{clustering\_unsupervised\_discovery/k\_means.py.jinja} 
    & K-means \\
\bottomrule
\end{tabular}
\end{table}

Templates are written manually and therefore follow a simple 
structure. Adding a template for a new method only requires 
creating a template file with the correct method 
name in the appropriate family folder. This makes the 
template system straightforward to extend.